% Generated draft: requires mathematical and source review.
\documentclass[11pt]{article}
\usepackage[margin=25mm]{geometry}
\usepackage{fontspec}
\setmainfont{Latin Modern Roman}
\usepackage{amsmath,amssymb,mathtools,unicode-math}
\setmathfont{Latin Modern Math}
\usepackage{enumitem,needspace,xurl,hyperref}
\setlength{\parindent}{0pt}
\setlength{\parskip}{5pt}
\setlength{\emergencystretch}{3em}
\newcommand{\N}{\mathbb{N}}
\newcommand{\Z}{\mathbb{Z}}
\newcommand{\Q}{\mathbb{Q}}
\newcommand{\R}{\mathbb{R}}
\newcommand{\C}{\mathbb{C}}
\newcommand{\F}{\mathbb{F}}
\newcommand{\A}{\mathbb{A}}
\newcommand{\PP}{\mathbb{P}}
\newcommand{\E}{\mathbb{E}}
\newcommand{\eps}{\varepsilon}
\newcommand{\dd}{\,\mathrm{d}}
\begin{document}
\section*{UnsolvedMath ID: 701181}
\textbf{Source problem code:} catalogue entry
\subsection{Definitions and mathematical statement}
Let \(p>2\) be a prime. A pro-\(p\) group is a topological group that is an inverse limit of finite \(p\)-groups. Automorphisms of pro-\(p\) groups are understood here to be continuous. For an automorphism \(\varphi\) of a pro-\(p\) group \(G\), define its fixed-point subgroup by
\[
C_G(\varphi)=\{g\in G:\varphi(g)=g\}.
\]
For a positive integer \(d\), let \(F_d^{(p)}\) be the free pro-\(p\) group on generators \(x_1,\ldots,x_d\). Given \(w\in F_d^{(p)}\) and \((g_1,\ldots,g_d)\in G^d\), define \(w(g_1,\ldots,g_d)\) as the image of \(w\) under the unique continuous homomorphism sending each \(x_i\) to \(g_i\). In this draft, saying that \(G\) satisfies a pro-\(p\) identity means that
\[
\exists d\geq 1\ \exists w\in F_d^{(p)}\setminus\{1\}\ \forall(g_1,\ldots,g_d)\in G^d,
\qquad w(g_1,\ldots,g_d)=1.
\]
The source does not explicitly define ``pro-\(p\) identity''; this finite-variable interpretation requires confirmation against the original source. The word \(w\) is allowed to depend on \(G\), and need not be an ordinary abstract group word.

The question is whether, for every prime \(p>2\), every pro-\(p\) group \(G\), and every continuous automorphism \(\varphi\) of exact order \(2\), the hypothesis that \(C_G(\varphi)\) is abelian implies that \(G\) satisfies a pro-\(p\) identity in the sense above. No finite-generation hypothesis is supplied.

The supplied background says that an affirmative answer would generalize a result attributed to A. N. Zubkov asserting that nonabelian free pro-\(p\) groups cannot be represented by \(2\times2\) matrices. The coefficient setting and precise meaning of representation in that background assertion are not specified, so it is retained as attributed context rather than a fully formulated auxiliary theorem.
\subsection{Short English statement}
For an odd prime \(p\), must every pro-\(p\) group admitting an involutory automorphism with abelian fixed-point subgroup satisfy a nontrivial pro-\(p\) identity? The draft interprets this as a nonidentity element of a free pro-\(p\) group on finitely many generators that evaluates to the identity on every tuple in \(G\).
\subsection{Sources}
\begin{itemize}
\item Local UnsolvedMath registry, permanent ID \(701181\), exact source problem code \texttt{catalogue entry}, title \emph{Kourovka Notebook Problem 15.93}. The supplied record labels the problem open; this label is supplied curation, not an independently verified assessment.
\item \emph{The Kourovka Notebook (Unsolved Problems in Group Theory)}, original problem 15.93, 15th issue (2002), proposed by A. Jaikin-Zapirain. The record also identifies the 21st issue edition (2026), arXiv:1401.0300v48. The collection description names E. I. Khukhro and V. D. Mazurov as editors. These bibliographic details are supplied and have not been independently checked.
\item A. N. Zubkov, \emph{Siberian Math. J.}, volume 28, no. 5 (1987), pp. 742--747, cited in the supplied statement for the assertion concerning nonabelian free pro-\(p\) groups and \(2\times2\) matrices. No article title or coefficient setting is supplied.
\end{itemize}
\end{document}
